\documentclass[11pt,a4paper]{article}

% ── Packages ──
\usepackage[utf8]{inputenc}
\usepackage[T1]{fontenc}
\usepackage{lmodern}
\usepackage[margin=1in]{geometry}
\usepackage{amsmath,amssymb}
\usepackage{graphicx}
\usepackage{booktabs}
\usepackage{hyperref}
\usepackage{listings}
\usepackage{xcolor}
\usepackage{natbib}
\usepackage{abstract}

% ── Listings style ──
\definecolor{codebg}{RGB}{248,248,248}
\definecolor{codeframe}{RGB}{200,200,200}
\lstset{
  basicstyle=\ttfamily\scriptsize,
  backgroundcolor=\color{codebg},
  frame=single,
  rulecolor=\color{codeframe},
  breaklines=true,
  breakatwhitespace=false,
  columns=fullflexible,
  keepspaces=true,
  showstringspaces=false,
  captionpos=b,
  aboveskip=8pt,
  belowskip=8pt,
  xleftmargin=3pt,
  xrightmargin=3pt,
  extendedchars=true,
  inputencoding=utf8,
}

\hypersetup{colorlinks=true, linkcolor=blue, citecolor=blue, urlcolor=blue}

% ── Title ──
\title{Emergent Multi-Agent Collaboration:\\Heterogeneous Cross-Model (Claude--Codex) \\ {\large Trial 3 Technical Report}}
\author{claude-and-codex research project}
\date{March 2026}

\begin{document}
\maketitle

% ── Abstract ──
\begin{abstract}
We study emergent collaboration between AI coding agents given no human direction.
In this trial, 2 agent(s) (Claude, Codex)
were launched simultaneously into a shared filesystem workspace with a minimal prompt
containing only identity, workspace path, and the other agent's name---no task, no
instructions, no time limit. The agents autonomously produced 184 lines of Python code across 2 file(s), with 12 passing tests. The resulting project was: Unknown Project.
This report documents the complete experimental procedure, inter-agent communication
transcripts, produced artifacts, and analysis of collaboration dynamics.
\end{abstract}

% ══════════════════════════════════════════════════════════════════════════════
\section{Introduction}

Recent work has demonstrated that large language model (LLM) agents can autonomously
collaborate when given access to a shared filesystem and minimal instructions
\citep{papailiopoulos2026}. In the ``when-claudes-meet'' experiment, two Claude Code
instances independently invented filesystem messaging protocols and built a complete
programming language in 12 minutes with zero human intervention.

This report presents one trial in a systematic study extending these findings across
three experimental configurations:
\begin{enumerate}
  \item \textbf{CC}: Homogeneous same-model (Claude + Claude)
  \item \textbf{CX}: Heterogeneous cross-model (Claude + Codex)
  \item \textbf{DCX}: Supervised heterogeneous (Director + Claude + Codex)
\end{enumerate}

This document covers the \textbf{CX} configuration, trial 3.

% ══════════════════════════════════════════════════════════════════════════════
\section{Related Work}

\citet{papailiopoulos2026} demonstrated that two Claude Code instances, given a shared
directory and the instruction ``find each other and build something together,''
independently converge on the same filesystem messaging protocol
(\texttt{hello} $\to$ \texttt{ack} $\to$ \texttt{proposal} $\to$ \texttt{build} $\to$
\texttt{done}). Their experiments produced a programming language (``Duo,'' 2{,}495 LOC,
41 tests) and a Battleship game.

Anthropic's Agent Teams feature \citep{anthropic2026teams} formalizes multi-agent
coordination with a lead-teammate architecture and shared task lists. Unlike emergent
collaboration, Agent Teams prescribes the coordination protocol.

Our work differs in two ways: (a) we compare same-model vs.\ cross-model vs.\
supervised configurations, and (b) we use truly minimal prompts with no prescribed task
or communication protocol.

% ══════════════════════════════════════════════════════════════════════════════
\section{Methodology}

\subsection{Experimental Design}

Each trial launches 2 agent process(es) simultaneously as
parallel subprocesses sharing a single filesystem directory. Agents communicate
exclusively through file creation and reading. No other channel is available. The
experiment runs with a 10-minute timeout.

\subsection{Agent Infrastructure}

Claude agents: \texttt{claude -p --dangerously-skip-permissions}
(unrestricted filesystem access and tool usage).

Codex agents: \texttt{codex exec -C <workspace> -s danger-full-access --skip-git-repo-check}
(equivalent filesystem permissions; the \texttt{-s danger-full-access} flag was required
after discovering that the default \texttt{--full-auto} sandbox blocks writes to shared
directories).

\subsection{Prompt Design}

The prompt contains \emph{only} identity and awareness---no task, no instructions:

\noindent\textbf{Claude:}
\begin{lstlisting}
You are "Claude". Your shared workspace is: <path>
The other agent is "Codex". You can communicate through files in the workspace. This workspace is yours -- you are free to use it however you want.
\end{lstlisting}

\noindent\textbf{Codex:}
\begin{lstlisting}
You are "Codex". Your shared workspace is: <path>
The other agent is "Claude". You can communicate through files in the workspace. This workspace is yours -- you are free to use it however you want.
\end{lstlisting}

\subsection{Metrics}

We measure: (1)~project chosen, (2)~lines of code produced, (3)~number of files,
(4)~communication files exchanged, (5)~test presence and pass rate,
(6)~whether both agents contributed code, (7)~collaboration patterns and failure modes.

% ══════════════════════════════════════════════════════════════════════════════
\section{Experimental Setup}

\begin{table}[h]
\centering
\begin{tabular}{ll}
\toprule
\textbf{Parameter} & \textbf{Value} \\
\midrule
Setting & Heterogeneous Cross-Model (Claude--Codex) \\
Trial & 3 \\
Agents & Claude, Codex \\
Timeout & 10 minutes \\
Human direction & None \\
\bottomrule
\end{tabular}
\caption{Experiment configuration.}
\label{tab:config}
\end{table}

% ══════════════════════════════════════════════════════════════════════════════
\section{Results}

\subsection{Project Output}

\begin{itemize}
  \item \textbf{Project}: Unknown Project
  \item \textbf{Python files}: 2
  \item \textbf{Lines of code}: 184
  \item \textbf{Communication files}: 1
  \item \textbf{Tests}: 12 passed
\end{itemize}

\subsection{Communication Protocol}

The agents exchanged 1 markdown file(s):
\begin{itemize}
  \item \texttt{HELLO\_CODEX.md}
\end{itemize}

% ══════════════════════════════════════════════════════════════════════════════
\section{Analysis \& Discussion}

\subsection{Collaboration Dynamics}

In the heterogeneous configuration, Claude (Anthropic) and Codex (OpenAI) exhibit complementary collaboration patterns. Claude typically proposes architecture with explicit interface contracts, while Codex tends to accept the proposal and focuses on polished user-facing components. Cross-model agents produce cleaner module boundaries.

\subsection{Observed Patterns}

\begin{enumerate}
  \item \textbf{Build-first strategy}: Claude consistently begins implementation before
    receiving explicit agreement, relying on well-structured interface contracts to
    minimize integration risk.
  \item \textbf{Universal messaging protocol}: All agents independently invent the same
    communication pattern (markdown files with structured proposals), confirming
    \citet{papailiopoulos2026}'s finding.
  \item \textbf{Graceful degradation}: When one agent is slow to respond, the other
    proceeds independently, designing code with fallback mechanisms.
\end{enumerate}

\subsection{Failure Modes}

\begin{enumerate}
  \item \textbf{File overwrite conflicts}: Agents may overwrite each other's files.
  \item \textbf{Proposal collision}: Both agents sometimes propose different projects
    simultaneously.
  \item \textbf{Passive waiting}: With minimal prompts, agents may default to waiting
    for instructions rather than taking initiative.
\end{enumerate}

% ══════════════════════════════════════════════════════════════════════════════
\section{Conclusion}

This trial demonstrates that 2 AI agent(s) can autonomously
produce working software (184 LOC, 12 passing tests) with zero human direction. The CX configuration yields clean module boundaries through cross-model complementarity.

Future work should investigate more complex tasks, conflict resolution protocols,
and the effect of prompt phrasing on agent activation.

% ══════════════════════════════════════════════════════════════════════════════
\bibliographystyle{plainnat}
\begin{thebibliography}{3}

\bibitem[Papailiopoulos(2026)]{papailiopoulos2026}
Dimitris Papailiopoulos.
\newblock When Claudes Meet: Emergent Collaboration Between AI Coding Agents.
\newblock GitHub: \texttt{anadim/when-claudes-meet}, 2026.

\bibitem[Anthropic(2026)]{anthropic2026teams}
Anthropic.
\newblock Orchestrate Teams of Claude Code Sessions.
\newblock Claude Code Documentation, 2026.

\bibitem[OpenAI(2026)]{openai2026codex}
OpenAI.
\newblock Codex CLI: Full-Auto Mode and Sandbox Policies.
\newblock Codex CLI Documentation, 2026.

\end{thebibliography}

% ══════════════════════════════════════════════════════════════════════════════
\appendix
\section{Communication Transcripts}

\subsection*{\texttt{HELLO\_CODEX.md}}
\begin{lstlisting}
# Hello Codex!

I'm Claude. This is our shared workspace (round 3).

We don't have a specific task from the supervisor yet. Let's pick something ourselves and build it.

## Proposal: Langton's Ant

A classic cellular automaton -- simple rules, emergent complexity. Perfect for a quick collaborative build.

**Rules:**
- Ant on a white cell: turn 90 right, flip cell to black, move forward
- Ant on a black cell: turn 90 left, flip cell to white, move forward

**Plan:**
1. `ant.py` -- Core simulation (grid + ant logic)
2. `test_ant.py` -- Tests for correctness
3. Terminal renderer that prints the grid state

If you have a different idea, drop it in a file and I'll pivot. Otherwise, I'll start on the core simulation.

-- Claude
\end{lstlisting}

\section{Source Code}

Total: 2 Python files, 184 lines.

\subsection*{\texttt{ant.py}}
\begin{lstlisting}[language=Python]
"""Langton's Ant -- cellular automaton simulation."""

from __future__ import annotations

# Directions: 0=up, 1=right, 2=down, 3=left
DIRECTIONS = [(0, -1), (1, 0), (0, 1), (-1, 0)]  # (dx, dy)


class LangtonAnt:
    """Simulates Langton's Ant on an infinite grid."""

    def __init__(self, x: int = 0, y: int = 0, direction: int = 0):
        self.x = x
        self.y = y
        self.direction = direction  # 0=up, 1=right, 2=down, 3=left
        self.black_cells: set[tuple[int, int]] = set()
        self.step_count = 0

    def step(self) -> None:
        """Execute one step of the simulation."""
        pos = (self.x, self.y)
        if pos in self.black_cells:
            # On black: turn left, flip to white
            self.direction = (self.direction - 1) % 4
            self.black_cells.discard(pos)
        else:
            # On white: turn right, flip to black
            self.direction = (self.direction + 1) % 4
            self.black_cells.add(pos)
        # Move forward
        dx, dy = DIRECTIONS[self.direction]
        self.x += dx
        self.y += dy
        self.step_count += 1

    def run(self, steps: int) -> None:
        """Run the simulation for a given number of steps."""
        for _ in range(steps):
            self.step()

... (85 lines total)
\end{lstlisting}

\subsection*{\texttt{test\_ant.py}}
\begin{lstlisting}[language=Python]
"""Tests for Langton's Ant simulation."""

from ant import LangtonAnt, DIRECTIONS


def test_initial_state():
    ant = LangtonAnt()
    assert ant.x == 0
    assert ant.y == 0
    assert ant.direction == 0
    assert ant.step_count == 0
    assert len(ant.black_cells) == 0


def test_first_step_turns_right_on_white():
    ant = LangtonAnt()
    ant.step()
    # Was on white: turn right (0->1), flip to black, move right
    assert (0, 0) in ant.black_cells  # original cell now black
    assert ant.direction == 1  # facing right
    assert ant.x == 1
    assert ant.y == 0
    assert ant.step_count == 1


def test_second_step():
    ant = LangtonAnt()
    ant.step()  # white -> right, at (1,0) facing right
    ant.step()  # white -> right again, at (1,1) facing down
    assert (1, 0) in ant.black_cells
    assert ant.direction == 2  # facing down
    assert ant.x == 1
    assert ant.y == 1


def test_return_to_black_cell_turns_left():
    ant = LangtonAnt()
    # Run 4 steps: ant makes a small pattern and revisits area
    ant.run(4)
    # After 4 steps the ant should have visited and flipped cells
... (99 lines total)
\end{lstlisting}

\section{Test Results}

\begin{lstlisting}
============================= test session starts =============================
platform win32 -- Python 3.12.7, pytest-9.0.2, pluggy-1.5.0 -- C:\Users\hyogon.ryu\AppData\Local\miniconda3\python.exe
cachedir: .pytest_cache
rootdir: C:\Users\hyogon.ryu\Desktop\project\claude-and-codex
configfile: pyproject.toml
plugins: anyio-4.10.0, asyncio-1.3.0, cov-7.0.0
asyncio: mode=Mode.AUTO, debug=False, asyncio_default_fixture_loop_scope=None, asyncio_default_test_loop_scope=function
collecting ... collected 11 items

test_ant.py::test_initial_state PASSED                                   [  9%]
test_ant.py::test_first_step_turns_right_on_white PASSED                 [ 18%]
test_ant.py::test_second_step PASSED                                     [ 27%]
test_ant.py::test_return_to_black_cell_turns_left PASSED                 [ 36%]
test_ant.py::test_symmetry_after_104_steps PASSED                        [ 45%]
test_ant.py::test_highway_emerges PASSED                                 [ 54%]
test_ant.py::test_custom_start_position PASSED                           [ 63%]
test_ant.py::test_bounds PASSED                                          [ 72%]
test_ant.py::test_render_contains_ant PASSED                             [ 81%]
test_ant.py::test_render_contains_black_cells PASSED                     [ 90%]
test_ant.py::test_directions_are_four PASSED                             [100%]

============================= 11 passed in 0.03s ==============================
\end{lstlisting}

\end{document}
